\section{Method}
\label{sec:method}
\subsection{Deterministic BEV representation}
\label{sec:encoding}
The input is a LiDAR scan of coordinates $(x,y,z)$ and reflectance. We retain finite points inside $x\in[0,70.4]$ m, $y\in[-40,40]$ m, and $z\in[-2.5,1.0]$ m, with a small boundary exclusion used by the encoder. Horizontal cells measure $0.1\times0.1$ m, giving a channel-first tensor of size $14\times800\times704$. The fourteen-channel representation, termed Hist14, comprises four height-bin counts; maximum, mean, span, and standard deviation of height; maximum and mean reflectance; total density; occupancy; and two mean horizontal offsets from the cell center. Empty cells are zero. Height statistics use $u=(z+2.5)/3.5$, and the four histogram bins divide this normalized interval equally. Each count $n$ is mapped to $\min\{1,\log(1+n)/\log(33)\}$. Reflectance is clipped to $[0,1]$, and the offsets are divided by cell width and clipped to $[-0.5,0.5]$. This fixed encoding preserves vertical summaries without a learned point encoder. Supplementary Note~S1 gives the channel definitions and supplied channel visualization.

\subsection{Backbone and scale-gated fusion}
The stem applies two $3\times3$ convolutions with 32 channels, the first with stride two. Three stages then produce $C_3$, $C_4$, and $C_5$ at strides 4, 8, and 16, with 48, 96, and 128 channels and depths 3, 4, and 2. A stage block applies a $7\times7$ depthwise convolution, batch normalization, pointwise expansion by a factor of 2.5, SiLU activation, pointwise projection, and batch normalization. Its output is $X+\gamma\odot F(X)$, where the channel scale $\gamma$ is initialized to $10^{-5}$. This combines large spatial kernels with narrow learned channel transformations~\cite{liu2022convnext}. The evaluated main configurations do not use structural reparameterization.

The top-down neck projects $C_5$ and $C_4$ to 48 channels and $C_3$ to 32 channels. A refined, bilinearly upsampled feature $U_j$ is fused with a lateral feature $L_j$ through
\begin{equation}
 \begin{aligned}
 P_j&=U_j+G_j\odot L_j,\\
 G_j&=2\sigmoid\!\left(\operatorname{DWConv}_{3\times3}(L_j+U_j)\right).
 \end{aligned}
\end{equation}
Zero gate initialization gives $G_j=1$. The first refinement is depthwise; the next projection maps 48 channels to 32. A final $3\times3$ convolution produces the $32\times200\times176$ prediction feature. Historical comparisons also use fusion width 24 and prediction width 16, as specified in the supplement.

\subsection{Context and residual attention}
The optional focal context module at $C_4$ follows the modulation principle of FocalNet~\cite{yang2022focalnet}. A normalized input is projected into a 64-channel query, a 64-channel context seed, and four gates. Three sequential depthwise $3\times3$ layers with dilations 1, 2, and 3 generate local context levels; their final level also provides a global spatial mean. A softmax over the four gates weights the three local levels and global context. A pointwise context transform modulates the query elementwise, and a projection returns 96 channels through a residual scale initialized to $10^{-3}$. The softmax and geometry-sensitive arithmetic use FP32.

P16 additionally applies SimAM~\cite{yang2021simam} at $C_3$. For each channel, let $\mu$ be its spatial mean, $d=(X-\mu)^2$, and $v=\sum d/(HW-1)$. The attention transform and its residual application are
\begin{equation}
 \begin{aligned}
 T(X)&=X\odot\sigmoid\!\left(\frac{d}{4(v+\lambda)}+\frac12\right),\\
 X'&=X+\beta[T(X)-X],
 \end{aligned}
\end{equation}
with $\lambda=10^{-4}$ and a learned scalar $\beta$ initialized to $10^{-3}$. SimAM's gate itself has no learned parameters; this residual wrapper adds one. ECA and a stride-two detail path are evaluated as separate alternatives.

\subsection{Grouped BEV and three-dimensional heads}
\label{sec:heads}
Two task groups predict Car and Pedestrian/Cyclist. Each group has classification, center offset, log-size, doubled-yaw, and optional quality branches. Geometric outputs have two channels per branch and are shared across classes within a group. The yaw representation $(\cos 2\theta,\sin 2\theta)$ decodes as $\theta=\tfrac12\operatorname{atan2}(s,c)$, representing orientation modulo $\pi$. Centers are decoded from metric offsets relative to output cell centers, and width and length are exponentiated from log-size predictions. Gaussian heatmaps provide center supervision~\cite{zhou2019objects}; the radius uses overlap 0.1 and a minimum of four output cells. Regression targets use nearest-center assignment.

The archived BEV models use independent branches, each with two $3\times3$ convolutions and a pointwise output layer. In Final3D, each group shares the first $3\times3$ convolution while retaining a private second convolution and output layer per branch. An additional branch predicts LiDAR bottom height $z_b$ and $\log h$. Its masked Smooth-L1 objective, with coefficient one, is added outside the adaptive BEV task pool. The two group losses receive equal weights of 0.5. Calibrated bottom-center transformation and corner projection produce the camera-coordinate boxes required by the reference evaluator. The model predicts upright yaw-only cuboids and does not resolve front versus back; orientation similarity is therefore omitted.

\subsection{Oriented geometric alignment}
\label{sec:oga}
The baseline combines Gaussian focal classification with masked coordinate losses. OGA replaces the coordinate L1 terms with Smooth-L1 and adds a projection objective and a normalized corner term. For a predicted footprint $B_p$ and target $B_t$, we project their corners onto the length and width axes of both rectangles, following the multi-axis construction of MGIoU~\cite{le2026mgiou}. For axis $k$, let $I_k$, $U_k$, and $C_k$ denote projected intersection, union, and enclosing interval lengths. We compute
\begin{equation}
 g_k=\frac{I_k}{U_k}-\frac{C_k-U_k}{C_k},\qquad
 \bar g=\frac14\sum_{k=1}^{4}g_k,\qquad
 L_{\mathrm{proj}}=\frac{1-\bar g}{2}.
\end{equation}
Denominators use a small numerical floor. The corner objective takes the smaller mean L1 corner distance under canonical and $\pi$-equivalent corner ordering, divided by the target footprint diagonal. Its coefficient is one. The geometry package therefore measures both projected alignment and normalized corner displacement, while preserving the orientation symmetry of the head.

For $M$ BEV tasks, OGA uses bounded adaptive task weights
\begin{equation}
 b_i=3\tanh(s_i/3),\qquad
 w_i=M\frac{\exp(b_i/2)}{\sum_j\exp(b_j/2)}.
\end{equation}
An exponential moving average $\mu_i$ of each detached task loss uses momentum 0.99 and a floor of $10^{-4}$. The group objective is $\bar\mu\sum_i w_iL_i/\mu_i$, where $\bar\mu$ is the mean running scale. This bounded weighting and scale normalization are an implemented training heuristic; they do not imply equal gradient magnitudes or a probabilistic uncertainty model. Our ablation evaluates the package as a whole.

\subsection{Localization quality and inference}
\label{sec:quality}
IQA predicts a scalar quality $q$ per cell. Its positive-location binary cross-entropy target is $q^*=\sg[\operatorname{clip}(\bar g,0,1)]$, with coefficient one. Since $\bar g$ averages projection GIoUs, this target is a geometric surrogate rather than exact rotated IoU. The score used for candidate ranking is
\begin{equation}
 s=p_{\mathrm{cls}}^{1-\alpha}q^\alpha,\qquad \alpha=0.5.
\end{equation}
Configurations without IQA use classification scores. Per-class peak extraction, a score threshold of 0.05, classwise rotated BEV suppression at overlap 0.1, and a global maximum of 500 detections form the evaluated decoding policy. The vertical branch does not change the quality target or suppression domain. Q-OGA, a comparison inspired by quality focal learning~\cite{li2020gfl}, couples classification to the same geometric surrogate, uses softened corner assignment and range weighting, and has no separate IQA branch. Its full settings appear in Supplementary Note~S2.
